\chapter{Exercise 9}

% ---------------- Exercise 9.1 ----------------
\section{Exercise 9.1}
\subsection{Problem}
Let $A$ be declared as a $[2..4] \times [1..5]$ array of integers, and each integer occupies 4 bytes.
What is the translation result of input token string: $x := A[3,2]$?

\begin{itemize}
    \item Tips: use the translation scheme for Pascal.
\end{itemize}

\subsection{Solution}
For Pascal arrays, the row-major offset of $A[3,2]$ is
\[
(3 - 2) \times 5 + (2 - 1) = 6.
\]
Since each integer occupies 4 bytes, the byte offset is $6 \times 4 = 24$.
Thus the generated translation is
\begin{align*}
t_1 &= 3 - 2 \\
t_2 &= t_1 \times 5 \\
t_3 &= 2 - 1 \\
t_4 &= t_2 + t_3 \\
t_5 &= 4 \times t_4 \\
t_6 &= base_A + t_5 \\
x &= t_6
\end{align*}

% ---------------- Exercise 9.2 ----------------
\section{Exercise 9.2}
\subsection{Problem}
Let $a$ be declared as a $5 \times 6$ array of integers, and each integer occupies 4 bytes.
What is the translation result of input token string: $i = a[3][2]$?

\begin{itemize}
    \item Tips: use the translation scheme for C/C++.
\end{itemize}

\subsection{Solution}
In C/C++, arrays are stored in row-major order.
The byte offset of $a[3][2]$ is
\[
(3 \times 6 + 2) \times 4 = 80.
\]
So the translation result is
\begin{align*}
t_1 &= 3 \times 6 \\
t_2 &= t_1 + 2 \\
t_3 &= 4 \times t_2 \\
t_4 &= a[t_3] \\
i &= t_4
\end{align*}

% ---------------- Exercise 9.3 ----------------
\section{Exercise 9.3}
\subsection{Problem}
What is the translation result of input token string: \texttt{x < 100 || x > 200 \&\& x != y}?

\begin{itemize}
    \item Tips: use the translation scheme for boolean expressions with backpatching.
    \item Suppose that the start position of the generated code is 100.
\end{itemize}

\subsection{Solution}
Using short-circuit evaluation and backpatching, one valid generated code sequence is
\begin{align*}
100 &: \ \text{if } x < 100 \ \text{goto } \_ \\
101 &: \ \text{goto } 102 \\
102 &: \ \text{if } x > 200 \ \text{goto } 104 \\
103 &: \ \text{goto } \_ \\
104 &: \ \text{if } x \neq y \ \text{goto } \_ \\
105 &: \ \text{goto } \_
\end{align*}

The blank targets are filled by backpatching after the true and false lists of the whole expression are known.